\documentclass[../main.tex]{subfiles}
\begin{document}

\chapter{命令スケジューリング}
\lessoninfo{
  video={Lesson 7，動画 1--30},
  textbook={H\&P \cite{hennessy2017} §3.4--3.5；\cite{tomasulo1967}},
  keywords={アウトオブオーダ実行，Tomasulo のアルゴリズム，命令キュー，リザベーションステーション，レジスタエイリアステーブル，タグ，共通データバス，発行，ディスパッチ，結果書き込み，古くなった結果，最古優先},
}

第6章では，プログラムにはインオーダプロセッサが利用できるよりはるかに
多くの命令レベル並列性が含まれていること，そしてそれに近づくには，
レジスタをリネーミングし命令をアウトオブオーダに実行するプロセッサが
必要であることを示した．この章では，そのようなプロセッサが，どの命令を
実行するかをサイクルごとにどのように決めるかを述べる．その方法が
Tomasulo のアルゴリズムである．これはレジスタリネーミングと，オペランド
の揃った命令の検出とを組み合わせたもので，今日のアウトオブオーダ
プロセッサでも依然として基礎となっている．命令がたどる 3つのステップを
追い，複数の命令や結果が同じハードウェアを奪い合うときに生じる選択を
調べ，完全な例をトレースし，最後に各命令がいつ実行されるかを求める簡便法
を導く．

\section{IPC の向上}
\label{sec:l07-ipc}

実際のプログラムの ILP は~1 を大きく上回るので，IPC をそれに近づけよう
とするプロセッサは，命令間の 4 種類の依存に対処しなければならない．
\source{Lesson 7，動画 1--2；H\&P §3.1．}それぞれに対策がある．
\begin{description}
  \item[制御依存] 分岐に続く命令が分岐の実行前にパイプラインに入れるのは，
    分岐の結果を推測する場合だけである．分岐予測（第4章）と，予測の難しい
    分岐に対するプレディケーション（第5章）がこれに対処する．
  \item[偽のデータ依存] WAW 依存と WAR 依存は，命令がレジスタ名を再利用する
    というだけの理由で存在する．レジスタリネーミング（第6章）がこれらを
    除去する．
  \item[真のデータ依存] RAW 依存は除去できない．命令は，自分が読む値が計算
    されるまで待たなければならない．\caution{依存はプログラムの性質であり，
    ハザードはそれを解消しなければならないパイプラインの性質である．
    H\&P §3.1．}避けられるのは，その命令に依存しない
    命令までが一緒に待つことである．
  \item[構造依存] 命令はハードウェアも奪い合う．リザベーションステーション，
    実行ユニット，そして後述するバスである．対策は，これらの資源をより多く
    備えた幅の広いプロセッサ（第6章）であり，それによって同じサイクルに
    より多くの命令を発行し実行できる．
\end{description}

\begin{definition}[アウトオブオーダ実行]
  \label{def:l07-ooo}
  プログラム順で先行する命令がまだ待っている間でも，各命令を，それが読む
  値が得られ次第実行することを
  \term[あうとおぶおーだじっこう]{アウトオブオーダ実行}[out-of-order execution]%
  という．
\end{definition}

アウトオブオーダ実行はリネーミングと組み合わせて用いられる．リネーミング
がなければ，命令が追い越せるのは WAW 依存も WAR 依存も持たない先行命令
だけとなり，ILP の多くが失われる．残るのはスケジューリングの問題である．
プロセッサは毎サイクル，どの命令のオペランドがすべて揃っているかを判定し，
そのうちどれを実行するかを決めなければならない．

\section{Tomasulo のアルゴリズム}
\label{sec:l07-tomasulo}

アウトオブオーダプロセッサで用いられるスケジューリング手法は，IBM
System/360 Model~91 で導入された．\source{Lesson 7，動画 3；H\&P §3.4；
\cite{tomasulo1967}．}

\begin{definition}[Tomasulo のアルゴリズム]
  \label{def:l07-tomasulo}
  Tomasulo が 1967 年に発表したハードウェアアルゴリズムで，レジスタを
  リネーミングし，オペランドが揃った命令を判定することによって，各命令を
  そのオペランドが計算され次第実行させるものを
  \term[とますろのあるごりずむ]{Tomasulo のアルゴリズム}[Tomasulo's algorithm]%
  という．
\end{definition}

現在のプロセッサも本質的に同じ方法で命令をスケジューリング
する．\margin{現代の規模：Skylake はスケジューラに 97 命令，リオーダ
バッファ（第8章）に 224 命令を保持する．Intel 最適化マニュアル．}相違点を
\cref{tab:l07-modern} に示す．最初の 2つはアルゴリズムの適用範囲に
関わるだけだが，3つ目はその機構を変える．Tomasulo のアルゴリズムには
例外への備えがない．命令がアウトオブオーダに実行された後では，レジスタ
の内容は，例外を受け付けられるプログラム中のどの時点にも対応していると
は限らない．例外をサポートするには追加の構造が必要であり，それが第8章の
主題である．\sidenote{Model~91（1967）のクロックは \SI{60}{\nano\second} で，
このアルゴリズムを適用したのは浮動小数点演算だけであった．加算は 2 サイクル，
乗算は~3 サイクルであり，まさにこの理由でその割り込みは不正確であった．この
考え方はその後ほとんど使われず，1990 年代半ばに Pentium~Pro と PowerPC~604
で復活した．}

\begin{table}[htbp]
  \centering\small
  \caption{Tomasulo のアルゴリズムと現代のアウトオブオーダプロセッサ．}
  \label{tab:l07-modern}
  \begin{tabular}{@{}lll@{}}
    \toprule
    & Tomasulo のアルゴリズム & 現代のプロセッサ \\
    \midrule
    スケジューリングする命令 & 浮動小数点演算 & すべての命令 \\
    同時に処理中の命令       & 数個           & 数百個 \\
    例外                     & サポートしない & サポートする（第8章） \\
    \bottomrule
  \end{tabular}
\end{table}

\section{構成}
\label{sec:l07-organization}

\cref{fig:l07-tomasulo} は，Tomasulo のアルゴリズムを実行するハードウェア
を示す．\source{Lesson 7，動画 4；H\&P §3.4．}フェッチステージは，
フェッチした各命令を\term[めいれいきゅー]{命令キュー}[instruction queue]%
に追加する．命令キューは先入れ先出しのバッファであり，フェッチされた命令
を，リザベーションステーションへ移せるようになるまでプログラム順に保持
する．

\begin{figure}[htbp]
  \centering
  % Hardware organization of Tomasulo's algorithm: instruction queue, RAT and
% register file, reservation stations, execution units, common data bus, and
% the separate path of loads and stores (address adder, load and store
% buffers, memory).
\begin{tikzpicture}[x=1cm,y=1cm, font=\footnotesize\sffamily, >={Stealth[length=1.8mm]},
  blk/.style={draw, fill=ln-box, minimum height=0.6cm, inner sep=2pt, align=center},
  rs/.style={draw, fill=white, minimum width=1.5cm, minimum height=0.36cm, inner sep=0pt,
    font=\scriptsize\sffamily},
  qc/.style={draw, fill=white, minimum width=0.4cm, minimum height=0.5cm, inner sep=0pt},
  stp/.style={font=\scriptsize\bfseries\sffamily, text=ln-accent},
  bus/.style={line width=1.6pt, draw=ln-accent}]
  % The Japanese labels are wider than the English ones: shift the load and
  % store column to the right so that the "reservation stations" label fits
  % between the stations and the address adder (\lsd = 0 leaves EN unchanged).
  \iflangja{\def\lsd{0.8}}{\def\lsd{0}}
  % ---- fetch and instruction queue
  \node[blk, minimum width=1.1cm] (fetch) at (0.55,0) {\iflangja{フェッチ}{fetch}};
  \foreach \i in {0,...,5} { \node[qc] (q\i) at ({1.75+0.4*\i},0) {}; }
  \foreach \i in {2,...,5} { \node[qc, fill=ln-accent!15] at ({1.75+0.4*\i},0) {}; }
  \node[anchor=south, font=\scriptsize] at (2.75,0.28) {\iflangja{命令キュー}{instruction queue}};
  \draw[->] (fetch.east) -- (q0.west);
  % ---- issue: head of the queue to the stations, with the RAT
  \draw[->] (q5.east) -- (4.6,0) -- (4.6,-0.75);
  \node[stp, anchor=west] at (4.65,-0.3) {\iflangja{発行}{issue}};
  \draw (0.55,-0.75) -- ({9.3+\lsd},-0.75);
  \fill (4.6,-0.75) circle (1.3pt);
  % ---- RAT and register file
  \node[blk, fill=ln-accent!15, minimum width=1.1cm] (rat) at (0.55,-1.45) {RAT};
  \node[blk, minimum width=1.1cm, font=\scriptsize\sffamily] (regs) at (0.55,-2.35) {\iflangja{レジスタ\\ファイル}{register\\file}};
  \draw[<->] (rat.north) -- (0.55,-0.75);
  \draw (rat.south) -- (regs.north);
  % ---- reservation stations
  \foreach \n/\y in {1/-1.25,2/-1.61,3/-1.97} { \node[rs] (ad\n) at (2.9,\y) {AD\n}; }
  \foreach \n/\y in {1/-1.25,2/-1.61} { \node[rs] (ml\n) at (6.1,\y) {ML\n}; }
  \draw[->] (2.9,-0.75) -- (ad1.north);
  \draw[->] (6.1,-0.75) -- (ml1.north);
  \node[anchor=west, font=\scriptsize, text=ln-muted, align=left] at (6.95,-1.43)
    {\iflangja{リザベーション\\ステーション}{reservation\\stations}};
  % ---- load/store path: address adder, load and store buffers
  \node[blk, minimum width=1.6cm, minimum height=0.5cm] (aadd) at ({9.3+\lsd},-1.35) {\iflangja{アドレス加算器}{address adder}};
  \draw[->] ({9.3+\lsd},-0.75) -- (aadd.north);
  \node[anchor=south east, font=\scriptsize, text=ln-muted] at ({9.3+\lsd},-0.75) {LW, SW};
  \node[blk, minimum width=1.1cm, font=\scriptsize\sffamily] (sb) at ({8.7+\lsd},-2.3) {\iflangja{ストア\\バッファ}{store\\buffers}};
  \node[blk, minimum width=1.1cm, font=\scriptsize\sffamily] (lb) at ({9.9+\lsd},-2.3) {\iflangja{ロード\\バッファ}{load\\buffers}};
  \draw (aadd.south) -- ({9.3+\lsd},-1.8);
  \draw[->] ({9.3+\lsd},-1.8) -| (sb.north);
  \draw[->] ({9.3+\lsd},-1.8) -| (lb.north);
  % ---- execution units
  \node[blk, minimum width=1.5cm] (adder) at (2.9,-3.3) {\iflangja{加減算器}{adder}};
  \node[blk, minimum width=1.5cm] (mult) at (6.1,-3.3) {\iflangja{乗除算器}{multiplier}};
  \node[blk, minimum width=2.3cm] (mem) at ({9.3+\lsd},-3.3) {\iflangja{メモリ}{memory}};
  \draw[->] (ad3.south) -- (adder.north);
  \draw[->] (ml2.south) -- (mult.north);
  \draw[->] (sb.south) -- (sb.south |- mem.north);
  \draw[->] (lb.south) -- (lb.south |- mem.north);
  \node[stp, anchor=west] at (2.95,-2.6) {\iflangja{ディスパッチ}{dispatch}};
  \node[stp, anchor=west] at (6.15,-2.6) {\iflangja{ディスパッチ}{dispatch}};
  % ---- common data bus and its destinations
  \draw[bus] (0.55,-4.3) -- ({10.6+\lsd},-4.3);
  \node[anchor=north east, font=\scriptsize\bfseries\sffamily, text=ln-accent] at ({10.6+\lsd},-4.38)
    {CDB\enskip(\iflangja{結果書き込み}{write result})};
  \draw[->] (adder.south) -- (2.9,-4.3);
  \draw[->] (mult.south) -- (6.1,-4.3);
  \draw[->] (mem.south) -- ({9.3+\lsd},-4.3);
  \draw[->, draw=ln-accent] (0.55,-4.3) -- (regs.south);
  \draw[->, draw=ln-accent] (1.65,-4.3) -- (1.65,-1.61) -- ([xshift=-0.75cm]ad2.center);
  \draw[->, draw=ln-accent] (4.9,-4.3) -- (4.9,-1.43) -- ([xshift=-0.75cm]6.1,-1.43);
  \draw[->, draw=ln-accent] (1.65,-2.9) -- (1.25,-2.9) |- (rat.east);
  \draw[->, draw=ln-accent] (7.8,-4.3) -- (7.8,-2.3) -- (sb.west);
\end{tikzpicture}

  \caption{Tomasulo のアルゴリズムを用いるプロセッサの構成．発行は，RAT と
    レジスタファイルを用いて，キューの先頭の命令をリザベーションステーション
    に移す．ディスパッチは，準備のできた命令をその実行ユニットに送る．結果
    は共通データバスを通って，リザベーションステーション，RAT，レジスタ
    ファイル，ストアバッファに戻る．ロードとストアは，アドレス加算器と
    ロードバッファ・ストアバッファを経由する経路をとる．}
  \label{fig:l07-tomasulo}
\end{figure}

\begin{definition}[リザベーションステーション]
  \label{def:l07-rs}
  \term[りざべーしょんすてーしょん]{リザベーションステーション}[reservation station]%
  \index{RS|see{リザベーションステーション}}（RS）は，1つの命令を，それが
  命令キューを離れた時点から結果がブロードキャストされるまで保持する．
  RS は，使用中かどうか，演算の種類，各オペランドについてその値あるいは
  値がまだ計算されていなければその値のタグ，そして命令がすでに
  ディスパッチされたかどうかを記録する．
\end{definition}

リザベーションステーションは，それを担当する実行ユニットごとにまとめ
られる．この章の例では，1つの加減算器が担当する加算・減算用の 3つの
ステーション AD1，AD2，AD3 と，1つの乗除算器が担当する乗算・除算用の
2つのステーション ML1，ML2 を用いる．\margin{360/91 はまさにこの構成で
あった．浮動小数点加算器に 3つ，乗除算器に 2つのリザベーションステーションを
備えていた \cite{tomasulo1967}．}ステーションの名前は，その命令の
結果がブロードキャストされるときの\emph{タグ}である．

第6章のレジスタエイリアステーブルは，レジスタごとに 1つのエントリを持つ．
エントリが空であれば，レジスタファイルがそのレジスタの現在の値を保持して
いる．そうでなければ，エントリは，その値を生成する命令が入っている
ステーションのタグを保持する．したがってタグは，第6章の物理レジスタの
役割を果たす．すべての結果が新しい名前を受け取り，レジスタを読む命令は，
最も近い先行する書き込み命令の結果を参照する．

\begin{definition}[共通データバス]
  \label{def:l07-cdb}
  \term[きょうつうでーたばす]{共通データバス}[common data bus]%
  \index{CDB|see{共通データバス}}（CDB）は，実行ユニットが計算した各結果を，
  そのタグとともに，すべてのリザベーションステーション，RAT，レジスタ
  ファイルに運ぶ．
\end{definition}

ロードとストアは別の経路を用いる．加算器がメモリアドレスを計算した後，
ロードはアドレスを保持するロードバッファで，ストアはアドレスと格納する値
（CDB から受け取ることもある）を保持するストアバッファで待つ．ロードが
読んだ値は，他の結果と同様に CDB に載せられる（\cref{sec:l07-load-store}）．

すべての命令は，\cref{fig:l07-tomasulo} に示した 3つのステップを経る．
\begin{description}
  \item[発行] 命令キューの先頭の命令を，RAT を通じてオペランドと宛先を解決
    したうえで，空いているリザベーションステーションに置く．このステップを
    \term[はっこう]{発行}[issue]という．
  \item[ディスパッチ] オペランドがすべて揃った命令を，そのリザベーション
    ステーションから実行ユニットに送る．これを命令の
    \term[でぃすぱっち]{ディスパッチ}[dispatch]という．
  \item[結果書き込み] 結果を CDB に載せ，そこから待機中のリザベーション
    ステーション，RAT，レジスタファイルに届ける．このステップを
    \term[けっかかきこみ]{結果書き込み}[write result]，あるいは
    ブロードキャストという．\index{ぶろーどきゃすと@ブロードキャスト|see{結果書き込み}}
\end{description}

\section{発行}
\label{sec:l07-issue}

発行は，命令キューの先頭の命令を取り出し，次のステップを実行する．
\source{Lesson 7，動画 5--7；H\&P §3.5．}この章の例では 1 サイクルに
1 命令を発行する．幅の広いプロセッサは 1 サイクルに複数の命令を，やはり
プログラム順に発行する．
\begin{enumerate}
  \item プログラム順で次の命令を，キューの先頭から取り出す．
  \item 各オペランドの出所を決める．オペランドのレジスタの RAT エントリが
    空であれば，値はレジスタファイルにある．そうでなければ，エントリがその
    値のタグを示す．
  \item その演算用の空いているリザベーションステーションを探す．その種類の
    ステーションがすべて使用中であれば，命令はこのサイクルには発行されない．
    命令はキューの先頭に留まり，その後ろの命令も待つ．
  \item 命令をステーションに置く．演算の種類と，各オペランドについて
    レジスタファイルの値あるいは RAT のタグを格納する．
  \item 宛先レジスタの RAT エントリを，そのステーションの名前に設定する．
\end{enumerate}

命令は厳密にプログラム順に発行される．そのため RAT は，発行中の命令の
位置から見えるレジスタ値を表し，各オペランドは最も近い先行する書き込み
命令を参照する．リネーミングと同様に，オペランドの参照は宛先のエントリを
上書きする前に行う．\caution{発行はオペランドを待つことはなく，空いている
リザベーションステーションだけを待つ．}

\begin{example}
  \label{ex:l07-program}
  次のプログラムを，この章の実行例とする．
  \begin{lstlisting}[language={[x86masm]Assembler}, numbers=none]
I1:  MUL R1, R2, R3     ; R1 <- R2 * R3
I2:  ADD R4, R1, R5     ; R4 <- R1 + R5
I3:  SUB R6, R2, R5     ; R6 <- R2 - R5
I4:  ADD R1, R5, R3     ; R1 <- R5 + R3
I5:  DIV R2, R6, R3     ; R2 <- R6 / R3
I6:  MUL R3, R4, R1     ; R3 <- R4 * R1
I7:  SUB R5, R3, R2     ; R5 <- R3 - R2
\end{lstlisting}
  初期状態では RAT は空で，レジスタファイルは R1 $=1$，R2 $=14$，
  R3 $=4$，R4 $=3$，R5 $=6$，R6 $=7$ を保持する．サイクル~1 で，I1 が
  オペランド値 14 と 4 とともに ML1 に発行され，R1 の RAT エントリは ML1
  になる．サイクル~2 で，I2 が AD1 に発行される（\cref{fig:l07-issue}）．
  R1 の RAT エントリは ML1 を保持しているので，第 1 オペランドはタグ ML1
  となる．R5 のエントリは空なので，第 2 オペランドは値~6 となる．R4 の
  RAT エントリは AD1 になる．I3 と I4 がサイクル 3 と~4 で続く．I4 は再び
  R1 に書き込み，R1 のエントリの ML1 を AD3 で置き換える．これ以降に発行
  される命令のうち R1 を読むものはすべて，I1 ではなく I4 の結果を使う．
\end{example}

\begin{figure}[htbp]
  \centering
  % Issuing I2 (ADD R4, R1, R5) in cycle 2 of the running example: the RAT
% entry of R1 holds the tag ML1, the entry of R5 is empty so the value comes
% from the register file, and the RAT entry of R4 is set to AD1.
\begin{tikzpicture}[x=1cm,y=1cm, font=\footnotesize\sffamily, >={Stealth[length=1.8mm]},
  cl/.style={draw, fill=white, minimum width=0.9cm, minimum height=0.45cm, inner sep=0pt},
  nm/.style={cl, fill=ln-box, minimum width=0.8cm},
  new/.style={cl, fill=ln-accent!15},
  tg/.style={font=\footnotesize\itshape\sffamily},
  lab/.style={font=\scriptsize\sffamily, text=ln-muted},
  arr/.style={->, thick, ln-accent}]
  % ---- register file and RAT, one column per register
  \node[anchor=west] at (-0.8,1.05) {I2:\enskip\texttt{ADD R4, R1, R5}};
  \foreach \r/\x in {1/1.45,2/2.35,3/3.25,4/4.15,5/5.05,6/5.95} {
    \node[lab] at (\x,0.4) {R\r};
  }
  \node[lab] at (1.45,0.66) {src$_1$};
  \node[lab] at (4.15,0.66) {dst};
  \node[lab] at (5.05,0.66) {src$_2$};
  \node[lab, anchor=east] at (0.95,0) {\iflangja{レジスタ}{register file}};
  \node[lab, anchor=east] at (0.95,-0.45) {RAT};
  \foreach \v/\x in {1/1.45,14/2.35,4/3.25,3/4.15,6/5.05,7/5.95} {
    \node[cl] at (\x,0) {\v};
  }
  \node[cl, tg] (r1) at (1.45,-0.45) {ML1};
  \foreach \x in {2.35,3.25,5.05,5.95} { \node[cl] at (\x,-0.45) {}; }
  \node[new, tg] (r4) at (4.15,-0.45) {AD1};
  \node[cl, draw=none, fill=none] (r5) at (5.05,-0.45) {};
  % ---- reservation stations: adder group left, multiplier group right
  \foreach \n/\y in {1/-1.6,2/-2.05,3/-2.5} {
    \node[nm] (ad\n) at (0.4,\y) {AD\n};
    \node[cl, minimum width=0.8cm] (ado\n) at (1.2,\y) {};
    \node[cl] (ada\n) at (2.05,\y) {};
    \node[cl] (adb\n) at (2.95,\y) {};
  }
  \foreach \n/\y in {1/-1.6,2/-2.05} {
    \node[nm] (ml\n) at (4.4,\y) {ML\n};
    \node[cl, minimum width=0.8cm] (mlo\n) at (5.2,\y) {};
    \node[cl] (mla\n) at (6.05,\y) {};
    \node[cl] (mlb\n) at (6.95,\y) {};
  }
  \node[new, minimum width=0.8cm] at (1.2,-1.6) {ADD};
  \node[new, tg] (opa) at (2.05,-1.6) {ML1};
  \node[new] (opb) at (2.95,-1.6) {6};
  \node at (5.2,-1.6) {MUL};
  \node at (6.05,-1.6) {14};
  \node at (6.95,-1.6) {4};
  \node[lab] at (1.7,-2.95) {\iflangja{加減算用}{for the adder}};
  \node[lab] at (5.7,-2.5) {\iflangja{乗除算用}{for the multiplier}};
  % ---- operand lookups
  \draw[arr] (r1.south) -- (opa.north);
  \draw[arr] (5.05,-0.225) -- (r5.south) -- (opb.north);
\end{tikzpicture}

  \caption{\cref{ex:l07-program} のサイクル~2 における I2 の発行．R1 の RAT
    エントリからタグ ML1 が得られ（タグは斜体），R5 のエントリは空なので
    値~6 がレジスタファイルから取られ，R4 のエントリは AD1 に設定される．
    網掛けのセルがこのステップで書き込まれる．}
  \label{fig:l07-issue}
\end{figure}

\section{ディスパッチ}
\label{sec:l07-dispatch}

命令は，すべてのオペランドが値になるまでリザベーションステーションで
待つ．\source{Lesson 7，動画 8；H\&P §3.5．}CDB に結果が現れるたびに，
すべてのリザベーションステーションが自身のオペランドのタグをバス上の
タグと比較し，一致したところではタグをバスの値で置き換える．これを値の
\emph{取り込み}という．オペランドがすべて値になったステーションは，
\emph{準備のできた}命令を保持している．\caution{文献によって名前が逆に
なる．命令をステーションに置くことを\emph{ディスパッチ}，ユニットに送ることを
\emph{発行}と呼ぶ流儀もある．H\&P のステップは発行・実行・結果書き込みで
ある．}そのステーションを担当する実行
ユニットが空いていれば，準備のできた命令はディスパッチされる．すなわち，
その演算の種類とオペランド値がユニットに送られ，ユニットが実行を開始する．
命令の実行中もステーションは使用中のままである．ステーションの名前が，
結果をブロードキャストするときのタグだからである．

\subsection{準備のできた命令が複数ある場合}

実行ユニットが空いたときに複数の命令が準備できていることがあり，その中
から 1つを選ばなければならない．\source{Lesson 7，動画 9--10．}準備の
できた命令は必要な値をすべて持っているので，どれを選んでも正しい．選択が
影響するのは性能とハードウェアコストだけである．\sidenote{ブロードキャスト
されたタグを待機中の全オペランドと比較するウェイクアップと，ユニットごとに
準備のできた命令を 1つ選ぶセレクトは，あわせて 1 サイクルに収まらなければ
ならない．そうでなければ依存する命令を連続して実行できない．このループの
コストはステーション数とともに増えるので，リオーダバッファが数百命令を保持
するのに対し，スケジューラが保持するのは 100 程度である
\cite{palacharla1997}．}
\begin{description}
  \item[最古優先]\index{さいこゆうせん@最古優先}最初に発行された命令を
    ディスパッチする．他の条件が同じなら，古い命令ほど，処理中の命令のうち
    その後に発行されてその結果を待っている可能性のあるものが多い．したがって
    この方針は次の方針を近似し，しかも発行の順序を記録するだけで済む．
  \item[依存命令最多優先] 待機中の命令のうち最も多くがその結果を必要とする
    命令をディスパッチし，できるだけ多くの命令を早く準備完了にする．最良の
    性能が期待できるが，その命令を見つけるには，すべての候補のタグを
    すべてのリザベーションステーションのオペランドと比較する必要があり，
    遅すぎるうえにコストも大きすぎる．
  \item[ランダム] 選択論理が最も単純である．多くの命令が依存する命令を，
    重要でない命令の後ろで待たせることがあるので，性能が低下する．
\end{description}
最古優先は他の 2つの方針のよい折衷であり，この章の例で用いる規則である．
実行例では，サイクル~7 で加減算器が空いたとき，I2 と I4 がともに準備
できている（\cref{sec:l07-example}）．I4 のほうが 2 サイクル長く準備完了
の状態にあるが，最古優先により I2 がディスパッチされる．

\section{結果書き込み}
\label{sec:l07-write}

実行ユニットが命令の実行を終えると，結果書き込みのステップがその結果を
ブロードキャストする．\source{Lesson 7，動画 11；H\&P §3.5．}
\begin{enumerate}
  \item 結果と，命令のリザベーションステーションのタグを CDB に載せる．
  \item このタグを保持するすべてのオペランドが値を取り込む
    （\cref{sec:l07-dispatch}）．
  \item あるレジスタの RAT エントリがそのタグを保持していれば，値を
    レジスタファイルのそのレジスタに書き込み，RAT エントリを空にする．
    これにより，後で発行される命令はレジスタファイルから値を読む．
  \item リザベーションステーションを解放する．
\end{enumerate}
ステーションは宛先レジスタを記録する必要がない．タグを保持する RAT
エントリがそれを特定するからである．\cref{fig:l07-broadcast} は，実行例に
おける I2 の結果のブロードキャストを示す．

\begin{figure}[htbp]
  \centering
  % Write result of I2 in cycle 9 of the running example: tag AD1 and value 62
% on the CDB.  ML1 captures the value, the RAT entry of R4 matches, so R4 is
% written and its entry cleared, and AD1 is freed.
\begin{tikzpicture}[x=1cm,y=1cm, font=\footnotesize\sffamily, >={Stealth[length=1.8mm]},
  cl/.style={draw, fill=white, minimum width=0.9cm, minimum height=0.45cm, inner sep=0pt},
  nm/.style={cl, fill=ln-box, minimum width=0.8cm},
  new/.style={cl, fill=ln-accent!15},
  tg/.style={font=\footnotesize\itshape\sffamily},
  lab/.style={font=\scriptsize\sffamily, text=ln-muted},
  arr/.style={->, thick, ln-accent}]
  % ---- register file and RAT
  \foreach \r/\x in {1/1.45,2/2.35,3/3.25,4/4.15,5/5.05,6/5.95} {
    \node[lab] at (\x,0.4) {R\r};
  }
  \node[lab, anchor=east] at (0.95,0) {\iflangja{レジスタファイル}{register file}};
  \node[lab, anchor=east] at (0.95,-0.45) {RAT};
  \foreach \v/\x in {1/1.45,14/2.35,4/3.25,6/5.05,8/5.95} {
    \node[cl] at (\x,0) {\v};
  }
  \node[new] (g4) at (4.15,0) {62};
  \foreach \t/\x in {AD3/1.45,ML2/2.35,ML1/3.25,AD2/5.05} {
    \node[cl, tg] at (\x,-0.45) {\t};
  }
  \node[cl] at (5.95,-0.45) {};
  \node[new] (r4) at (4.15,-0.45) {};
  % ---- common data bus
  \draw[line width=1.6pt, ln-accent] (-0.1,-1.05) -- (7.5,-1.05);
  \node[anchor=west, font=\scriptsize\bfseries\sffamily, text=ln-accent, align=left]
    at (7.55,-1.05) {CDB\\\iflangja{タグ}{tag} \textit{AD1}, \iflangja{値}{value} 62};
  % ---- reservation stations
  \foreach \n/\y in {1/-1.6,2/-2.05,3/-2.5} {
    \node[nm] (ad\n) at (0.4,\y) {AD\n};
    \node[cl, minimum width=0.8cm] at (1.2,\y) {};
    \node[cl] at (2.05,\y) {};
    \node[cl] at (2.95,\y) {};
  }
  \foreach \n/\y in {1/-1.6,2/-2.05} {
    \node[nm] (ml\n) at (4.4,\y) {ML\n};
    \node[cl, minimum width=0.8cm] at (5.2,\y) {};
    \node[cl] at (6.05,\y) {};
    \node[cl] at (6.95,\y) {};
  }
  % AD1 (I2) is freed
  \node[lab] at (2.05,-1.6) {\iflangja{解放}{freed}};
  % AD2 (I7), AD3 (I4), ML2 (I5): no operand waits for AD1
  \node at (1.2,-2.05) {SUB};  \node[tg] at (2.05,-2.05) {ML1}; \node[tg] at (2.95,-2.05) {ML2};
  \node at (1.2,-2.5) {ADD};   \node at (2.05,-2.5) {6};         \node at (2.95,-2.5) {4};
  \node at (5.2,-2.05) {DIV};  \node at (6.05,-2.05) {8};         \node at (6.95,-2.05) {4};
  % ML1 (I6) captures the value in its first operand
  \node at (5.2,-1.6) {MUL};
  \node[new] (cap) at (6.05,-1.6) {62};
  \node[tg] at (6.95,-1.6) {AD3};
  % ---- effects of the broadcast
  \draw[arr] (4.15,-1.05) -- (g4.south);
  \draw[arr] (6.05,-1.05) -- (cap.north);
  \fill[ln-accent] (4.15,-1.05) circle (1.5pt) (6.05,-1.05) circle (1.5pt);
  \node[lab, anchor=west] at (4.2,-0.83) {\iflangja{一致}{match}};
\end{tikzpicture}

  \caption{実行例のサイクル~9 における I2 の結果書き込み．ML1（I6）はタグ
    AD1 の代わりに 62 を取り込む．R4 の RAT エントリは AD1 を保持している
    ので，R4 は 62 になり，エントリは空になる．AD1 は解放される．AD1 を
    保持するオペランドや RAT エントリは他にない．}
  \label{fig:l07-broadcast}
\end{figure}

\subsection{1 サイクルに複数の結果がある場合}
\label{sec:l07-cdb-conflict}

2つの実行ユニットが同じサイクルに実行を終えることがあるが，1本のバスは
一度に 1つの結果しか運べない．\source{Lesson 7，動画 12．}1つの解決法は，
ユニットごとに別々のバスを設けることである．その場合，すべての
リザベーションステーションのすべてのオペランドが自身のタグを各バスと
比較し，一致したバスから値を選ばなければならず，レジスタファイルには
バスごとに書き込みポートが必要になる．もともと大きい比較論理が，バスの
数に比例して増大する．もう1つの方法は，バスを 1本のままとし，一方の結果
を次のサイクルまで待たせることである．ここで用いる規則は，\emph{遅い}
ほうのユニットの結果を優先する．その命令はより前にディスパッチされた
ので古い命令である傾向があり，その結果を待っている命令も多い可能性が
高いからである．例では，結果がバスを待っているユニットは次の命令を開始
できないと仮定する．

\subsection{古くなった結果}
\label{sec:l07-stale}

結果がブロードキャストされる時点で，同じレジスタに書き込む後の命令が
すでに発行されていることがある．\source{Lesson 7，動画 13；H\&P §3.5．}
このときそのレジスタの RAT エントリは，バス上のタグではなく後の命令の
タグを保持しており，ステップ~3 は何の効果も持たない．レジスタファイルも
RAT も変化しない．結果はそのレジスタにとってのみ\emph{古くなった}%
\index{ふるくなったけっか@古くなった結果}ものであり，そこでそれを無視する
のは正しい．両方の命令の後では，レジスタは後の結果を保持していなければ
ならず，その結果は後の命令がブロードキャストするときに届く．先の結果も
なお必要だが，それを必要とするのは 2つの書き込み命令の間に発行された
命令だけである．それらはそのタグを保持しており，ステップ~2 で値を取り
込む．後の書き込み命令より後に発行された命令は，後のタグを参照する．

\begin{example}
  \label{ex:l07-stale}
  実行例では，I1 と I4 がともに R1 に書き込む．サイクル~6 で I1 がタグ
  ML1 のもとで 56 をブロードキャストするとき，I4 がサイクル~4 で発行
  されているので，R1 の RAT エントリはすでに AD3 を保持している．
  レジスタファイルの R1 は古い値~1 のままであり，エントリも AD3 のまま
  である．一方，I2 を保持する AD1 は~56 を取り込む．R1 が 10 になるのは，
  サイクル~11 で I4 が AD3 のもとでブロードキャストするときである．
\end{example}

\begin{keypoint}
  ブロードキャストは，そのタグを待つすべてのオペランドに値を届けるが，
  レジスタファイルと RAT を更新するのは，宛先の RAT エントリがまだその
  タグを保持している場合だけである．
\end{keypoint}

\section{1 サイクル内のステップ}
\label{sec:l07-cycle}

1つの命令から見ると，Tomasulo のアルゴリズムは，発行，リザベーション
ステーションでの待機と取り込み，ディスパッチ，実行，結果書き込みという
一連の流れである．\source{Lesson 7，動画 14--19；H\&P §3.5．}プロセッサ
全体から見ると，毎サイクルすべてのステップが，それぞれ異なる命令に対して
行われる．同じサイクルに，1つの命令が発行され，複数のステーションが
ブロードキャストされた値を取り込み，準備のできた命令が空いているユニット
にディスパッチされることがあり，そして 1つの結果がブロードキャストされる．

\emph{同じ}命令の 2つのステップが同じサイクルに行われうるかどうかは，
回路設計の問題である．最も単純な設計では，各ステップはサイクル開始時点の
構造の状態を調べ，すべての更新はサイクルの終わりに反映される．このとき
次のようになる．
\begin{itemize}
  \item サイクル $C$ で発行された命令がリザベーションステーションに入って
    いるのはサイクル $C+1$ からであり，ディスパッチされるのは早くても
    $C+1$ である．
  \item サイクル $C$ で最後に残っていたオペランドを取り込んだ命令は，
    サイクル $C+1$ から準備完了となり，ディスパッチされるのは早くても
    $C+1$ である．
  \item サイクル $C$ のブロードキャストで解放されたステーションは，
    サイクル $C+1$ に新しい命令を受け入れられる．
\end{itemize}
取り込みとディスパッチの間の遅延は，追加の論理，すなわち現在のサイクルに
バス上に到着するオペランドも考慮するディスパッチによって取り除くことが
でき，ブロードキャストで解放されたステーションも同様に同じサイクルで
再利用できる．追加の論理は，それが拡張するステップと同じサイクルに置かれる
ので，クロックサイクルを長くする．ただし，命令が発行されたサイクルに
ディスパッチされることはない．発行が命令をステーションに書き込んでいる
間，ディスパッチからはその命令がまだ見えないからである．

同じサイクルにおける発行と結果書き込みの間の 2つの相互作用には注意が
必要である．
\begin{itemize}
  \item 両者が同じレジスタの RAT エントリを変更することがある．古い書き
    込み命令のブロードキャストはエントリを空にし，新しい書き込み命令の
    発行はエントリを新しいタグに設定する．発行はプログラム順で後の命令に
    属するので，発行が優先されなければならない．言い換えれば，結果書き
    込みを先に，発行を後に適用する．そうしないと，後続の命令はレジスタ
    ファイルを読み，新たに発行された命令の結果を見落としてしまう．
  \item ある命令が，そのオペランドの 1つがブロードキャストされるサイクルに
    発行されることがある．最も単純な設計では，エントリはサイクルの終わり
    にならないと空にならないので，命令は RAT にタグを見つけるが，
    ステーションに入るのは値がバスを通過した後である．したがって発行は
    オペランドのタグをバスと比較し，値を直接取らなければならない．そう
    しないと，命令は決して来ないブロードキャストを待ち続ける．
    ブロードキャストによる更新が同じサイクルで発行から見えるなら，発行は
    単にレジスタファイルから新しい値を読む．
\end{itemize}

\begin{keypoint}
  最も単純な設計では，命令は 1 サイクルに高々 1 ステップしか進まない．
  サイクル $C$ で発行された命令がディスパッチされるのは早くても $C+1$
  であり，サイクル $W$ でブロードキャストされた結果を待つ命令が
  ディスパッチされるのは早くても $W+1$ である．
\end{keypoint}

\section{ロード命令とストア命令}
\label{sec:l07-load-store}

ロードとストアは，レジスタを介した依存と同様の，メモリを介した依存を
生む．\source{Lesson 7，動画 20；H\&P §3.5．}
\begin{lstlisting}[language={[x86masm]Assembler}, numbers=none]
SW  R1, 0(R2)      ; MEM[R2] <- R1
LW  R3, 8(R4)      ; R3 <- MEM[R4 + 8]
\end{lstlisting}
R4${}+8$ が R2 に等しければ，ロードは直前に格納された値を読む．これは
メモリを介した RAW 依存である．同じアドレスへの 2つのストアは WAW 依存を
持ち，先行するロードが読むアドレスへのストアは WAR 依存を持つ．RAT は
これらのうち偽の依存を除去できない．RAT がリネーミングするのはレジスタで
あってメモリ位置ではなく，また 2つのメモリアクセスが衝突するかどうかは，
両方のアドレスが計算されてはじめて分かるからである．\margin{Skylake は
72 個のロードと 56 個のストアを同時に処理中にできる．Intel 最適化
マニュアル．}

そのため Tomasulo のアルゴリズムは，ロードとストアをプログラム順に実行
する．ロードとストアは他の命令と同様に，オペランドを RAT を通じて解決して
発行され，アドレス加算器がアドレスを計算する．その後，ロードはロード
バッファで，ストアはストアバッファで待つ（\cref{fig:l07-tomasulo}）．
格納する値がまだ計算されていないストアは，それを CDB から取り込む．
メモリアクセスは，ロードとストアが発行された順に，これらのバッファから
行われる．\caution{プログラム順とはメモリアクセスの順序のことである．
アドレス自体はオペランドが揃い次第計算されるので，その計算はアウトオブ
オーダに行われる．}ロードは読んだ値を，演算の結果とまったく同じように，自身の
タグのもとで CDB にブロードキャストする．現代のプロセッサはロードと
ストアもアウトオブオーダに実行するが，そのためにはプロセッサがメモリを
介した依存を検出して処理する必要がある．これは第9章で扱う．\margin{第9章
では，アドレスがまだ分からないストアをロードが追い越すことを許し，2つの
アドレスが重なっていたかどうかを後から検査する．}

\section{完全な例}
\label{sec:l07-example}

ここで \cref{ex:l07-program} のプログラムを最後まで実行する．
\source{Lesson 7，動画 21--26；H\&P §3.5．}前提は，\cref{sec:l07-cycle}
の最も単純な設計に次の仮定を加えたものである．
\begin{itemize}
  \item 1 サイクルに発行する命令は高々 1つであり，必要な種類の空いている
    ステーションのうち番号の最も小さいものに発行する．
  \item 加減算器は ADD と SUB を 2 サイクルで，乗除算器は MUL と DIV を
    4 サイクルで実行する．\margin{参考までに：Skylake では浮動小数点の
    加算・乗算が 4 サイクル，倍精度除算が 13--14 サイクル，整数加算が~1
    サイクルである．Agner Fog, \emph{Instruction tables}．}サイクル $D$ で
    ディスパッチされ，実行時間が
    $k$ サイクルの命令は，サイクル $D$ から $D+k-1$ まで実行され，
    サイクル $D+k$ でブロードキャストする．そのサイクルには，ユニットは
    すでに次の命令を受け入れられる．\margin{ここでの各ユニットはパイプ
    ライン化されていない．サイクル $D$ から $D+k-1$ まで占有される．実際の
    加算器や乗算器は毎サイクル新しい演算を受け付ける．}
  \item CDB は 1本であり，乗除算器が加減算器より優先される
    （\cref{sec:l07-cdb-conflict}）．結果がバスを待っているユニットは新しい
    命令を受け入れない．準備のできた命令は最古優先でディスパッチされる．
\end{itemize}
\cref{tab:l07-trace} は各サイクルに起こることを列挙し，
\cref{tab:l07-state} はサイクル 5，6，7 の後のリザベーションステーション
と RAT を示す．

\begin{table}[htbp]
  \centering\small
  \caption{完全な例の各サイクルの事象．取り込みは，それを引き起こす
    ブロードキャストとともに示す．}
  \label{tab:l07-trace}
  \begin{tabular}{@{}c>{\raggedright\arraybackslash}p{0.14\linewidth}c>{\raggedright\arraybackslash}p{0.48\linewidth}@{}}
    \toprule
    サイクル & 発行 & ディスパッチ & 結果書き込みと備考 \\
    \midrule
    1     & I1 $\to$ ML1 &    & \\
    2     & I2 $\to$ AD1 & I1 & \\
    3     & I3 $\to$ AD2 &    & I2 は ML1 を待つ \\
    4     & I4 $\to$ AD3 & I3 & R1 の RAT エントリが AD3 になる \\
    5     & I5 $\to$ ML2 &    & 全ステーションが使用中；I4 は準備完了，
                                加減算器は使用中 \\
    6     &              &    & I1: ML1 $=56$，R1 には古くなった結果；
                                AD1 が取り込む；I3 は実行終了，バス待ち；
                                I6 用の空きステーションなし \\
    7     & I6 $\to$ ML1 & I2 & I3: AD2 $=8$，R6 $\leftarrow 8$；ML2 が
                                取り込む；I2 と I4 が準備完了，I2 が古い \\
    8     & I7 $\to$ AD2 & I5 & \\
    9     &              & I4 & I2: AD1 $=62$，R4 $\leftarrow 62$；ML1 が取り込む \\
    10    &              &    & \\
    11    &              &    & I4: AD3 $=10$，R1 $\leftarrow 10$；ML1 が取り込む \\
    12    &              & I6 & I5: ML2 $=2$，R2 $\leftarrow 2$；AD2 が取り込む \\
    13--15 &             &    & I6 を実行中 \\
    16    &              &    & I6: ML1 $=620$，R3 $\leftarrow 620$；AD2 が取り込む \\
    17    &              & I7 & \\
    18    &              &    & \\
    19    &              &    & I7: AD2 $=618$，R5 $\leftarrow 618$ \\
    \bottomrule
  \end{tabular}
\end{table}

\begin{table}[htbp]
  \centering\small
  \caption{サイクル 5，6，7 の終了時のリザベーションステーションと RAT．
    タグは斜体で示し，「実行中」はディスパッチ済みの命令を表す．空の RAT
    エントリはレジスタファイルの値で示す．}
  \label{tab:l07-state}
  \begin{tabular}{@{}llll@{}}
    \toprule
    & サイクル 5 の後 & サイクル 6 の後 & サイクル 7 の後 \\
    \midrule
    AD1 & I2: \textit{ML1}, 6       & I2: 56, 6              & I2: 56, 6（実行中） \\
    AD2 & I3: 14, 6（実行中）       & I3: 14, 6（実行中）    & 空き \\
    AD3 & I4: 6, 4                  & I4: 6, 4               & I4: 6, 4 \\
    ML1 & I1: 14, 4（実行中）       & 空き                   & I6: \textit{AD1}, \textit{AD3} \\
    ML2 & I5: \textit{AD2}, 4       & I5: \textit{AD2}, 4    & I5: 8, 4 \\
    \midrule
    R1  & \textit{AD3} & \textit{AD3} & \textit{AD3} \\
    R2  & \textit{ML2} & \textit{ML2} & \textit{ML2} \\
    R3  & 4            & 4            & \textit{ML1} \\
    R4  & \textit{AD1} & \textit{AD1} & \textit{AD1} \\
    R5  & 6            & 6            & 6 \\
    R6  & \textit{AD2} & \textit{AD2} & 8 \\
    \bottomrule
  \end{tabular}
\end{table}

\paragraph{サイクル 1--5．}1 サイクルに 1 命令ずつ発行される．I1 と I3 は
発行時にすべてのオペランドが揃っており，次のサイクルで I1 は乗除算器に，
I3 は加減算器にディスパッチされる．I4 はサイクル~5 から準備完了だが，
加減算器は I3 を実行中である．I2 はタグ ML1 を，I5 は AD2 を待つ．
サイクル~5 の後，すべてのステーションが使用中となる．

\paragraph{サイクル 6．}両ユニットはサイクル~5 で実行を終えている．バスは
遅いほうの乗除算器の結果，すなわち ML1 のもとでの 56 を運ぶ．これは R1
にとっては古くなった結果である（\cref{ex:l07-stale}）．AD1 がそれを取り
込み，ML1 は解放される．I3 の結果は加減算器で待つ．サイクルの開始時点で
乗除算用のステーションが 2つとも使用中なので，I6 は発行されない．

\paragraph{サイクル 7．}I3 が AD2 のもとで 8 をブロードキャストする．R6 の
RAT エントリは AD2 を保持しているので，R6 は~8 になり，エントリは空に
なる．ML2 は I5 のためにこの値を取り込み，AD2 は解放される．加減算器は
新しい命令を受け入れられ，I2 と I4 がともに準備できているので，最古優先に
より I2 がディスパッチされる．I6 は前のサイクルで解放された ML1 に発行
され，AD1 と AD3 の 2つのタグを受け取る．

\paragraph{サイクル 8--19．}I7 は AD2 に発行され，ML1 と ML2 を待つ．I5 は
サイクル~8 に乗除算器で実行を始め，I4 は I2 が加減算器を解放するサイクル~9
に加減算器で実行を始める．I6 は I2（サイクル~9）と I4（サイクル~11）の
結果，そして I5 がサイクル~12 に解放する乗除算器を必要とする．I7 は，
サイクル~16 にブロードキャストされる I6 の結果を必要とする．サイクル~19
の最後のブロードキャストの後，レジスタファイルは R1 $=10$，R2 $=2$，
R3 $=620$，R4 $=62$，R5 $=618$，R6 $=8$ を保持する．\margin{この例の IPC は
$7/19 \approx 0.37$ である．ステーションとユニットが無制限でも，
I1$\to$I2$\to$I6$\to$I7 の連鎖に 17 サイクルを要する．}これはプログラムを
順に実行して得られる値とまったく同じである．

\begin{keypoint}
  この例の命令は I1，I3，I2，I4，I5，I6，I7 の順にブロードキャストし，
  各命令はオペランド，空いているステーション，空いているユニットが許す
  限り早く実行される．すべてのオペランドは発行時にタグを通じて解決され，
  古くなった結果はレジスタファイルに届かないので，最終的なレジスタ値は
  逐次実行の場合と同じになる．
\end{keypoint}

\section{構造を追跡しないタイミングの算出}
\label{sec:l07-timing}

すべての構造をサイクルごとにトレースするのは手間がかかる．タイミングだけ
が問題であれば，各命令の発行，ディスパッチ，結果書き込みのサイクルを
求めれば十分である．\source{Lesson 7，動画 27--29；H\&P §3.5．}命令 $n$
のこれらのサイクルを $I_n$，$D_n$，$W_n$ で，その実行時間を $k_n$ で表す．
\cref{sec:l07-example} の仮定のもとで，次が成り立つ．
\begin{description}
  \item[発行] $I_n$ は，$I_{n-1}$ より後で，必要な種類のステーションが
    空いている最初のサイクルである．命令 $m$ の結果書き込みで解放された
    ステーションは，サイクル $W_m+1$ から空いている．
  \item[ディスパッチ] $D_n$ は，次式を満たし，かつ実行ユニットが空いて
    いる最初のサイクルである．
    \begin{equation}
      D_n \;\ge\; 1 + \max\Bigl(I_n,\; \max_{p \in P(n)} W_p\Bigr) ,
      \label{eq:l07-dispatch}
    \end{equation}
    ここで $P(n)$ は，$n$ のオペランドを生成する命令（最も近い先行する
    書き込み命令）の集合である．\tip{$P(n)$ は，$n$ が読む各レジスタについての，
    最も近い先行する書き込み命令からなる（第6章）．}ユニットが空いている
    とは，他の命令 $m$ に占有されておらず（$m$ はサイクル $D_m$ から
    $W_m - 1$ までユニットを占有する），準備のできたより古い命令にも
    取られていないことである．
  \item[結果書き込み] $W_n = D_n + k_n$ である．ただし，そのサイクルの
    バスが優先権を持つ結果に取られていれば，それより後になる．
\end{description}
ユニットやバスは，先に準備のできた\emph{後の}命令に取られることがあるので，
行を単純にプログラム順に埋めていくことはできない．\tip{各ステップが待つ資源
は 1つである．発行は空きステーション，ディスパッチはオペランドと空きユニット，
結果書き込みはバスを待つ．}ユニットとバスは，
すべての命令にわたってサイクル順に割り当てる．各サイクルで，空いている
ユニットはそのサイクルに準備のできている命令のうち最も古いものを取り，
バスは優先権を持つ結果を取る．まだ準備のできていない古い命令のために
ユニットが空いたまま待つことはない．

\begin{figure}[htbp]
  \centering
  % Timing of the running example: issue, waiting in the reservation station,
% execution, completed result waiting for the CDB, and write result.
\begin{tikzpicture}[x=0.52cm,y=0.5cm, font=\scriptsize\sffamily,
  c/.style={minimum width=0.52cm, minimum height=0.5cm, inner sep=0pt, draw=ln-rule},
  iss/.style={c, fill=white, draw=black},
  ex/.style={c, fill=ln-accent!25, draw=black},
  wt/.style={c, fill=ln-caution!20, draw=black},
  wr/.style={c, fill=ln-accent, draw=black, text=white, font=\scriptsize\bfseries\sffamily},
  q/.style={c, fill=ln-muted!30, draw=black}]
  \foreach \c in {1,...,19} { \node[text=ln-muted] at (\c,1) {\c}; }
  \node[anchor=east, text=ln-muted] at (0.4,1) {\iflangja{サイクル}{cycle}};
  \foreach \c in {1,...,19} { \draw[ln-rule, very thin] (\c-0.5,0.6) -- (\c-0.5,-6.5); }
  \draw[ln-rule, very thin] (19.5,0.6) -- (19.5,-6.5);
  % name, issue, first execution cycle, last execution cycle, write result
  \foreach \n/\op/\i/\ea/\eb/\w [count=\r from 0] in
    {1/MUL/1/2/5/6, 2/ADD/2/7/8/9, 3/SUB/3/4/5/7, 4/ADD/4/9/10/11,
     5/DIV/5/8/11/12, 6/MUL/7/12/15/16, 7/SUB/8/17/18/19} {
    \node[anchor=east] at (0.4,-\r) {I\n\enskip\texttt{\op}};
    \node[iss] at (\i,-\r) {I};
    \pgfmathtruncatemacro{\wa}{\i+1}
    \pgfmathtruncatemacro{\wb}{\ea-1}
    \ifnum\wa<\ea
      \draw[ln-muted, thick] (\wa-0.5,-\r) -- (\wb+0.5,-\r);
    \fi
    \foreach \e in {\ea,...,\eb} { \node[ex] at (\e,-\r) {}; }
    \node[wr] at (\w,-\r) {W};
  }
  % I3: result complete in cycle 5, bus taken by I1 in cycle 6
  \node[wt] at (6,-2) {};
  % I6: no free multiplier station in cycle 6
  \node[q] at (6,-5) {};
  % ---- legend
  \begin{scope}[shift={(0,-7.6)}]
    \node[iss] at (1,0) {I}; \node[anchor=west] at (1.5,0) {\iflangja{発行}{issue}};
    \draw[ln-muted, thick] (3.9,0) -- (4.9,0);
    \node[anchor=west] at (5,0) {\iflangja{RS で待機}{waiting in RS}};
    \node[ex] at (9.4,0) {}; \node[anchor=west] at (9.9,0) {\iflangja{実行}{execution}};
    \node[wr] at (13.3,0) {W}; \node[anchor=west] at (13.8,0) {\iflangja{結果書き込み}{write result}};
  \end{scope}
  \begin{scope}[shift={(0,-8.7)}]
    \node[wt] at (1,0) {};
    \node[anchor=west] at (1.5,0) {\iflangja{実行済み，CDB 待ち}{done, waiting for the CDB}};
    \node[q] at (9.4,0) {};
    \node[anchor=west] at (9.9,0) {\iflangja{空き RS なし（命令キュー）}{no free RS (in queue)}};
  \end{scope}
\end{tikzpicture}

  \caption{実行例のタイミング．アウトオブオーダ実行により，I2 が待っている
    間に I3 が実行され，I4 が待っている間に I5 が実行を始める．I3（バス），
    I4（加減算器），I6（リザベーションステーション）の遅延は，
    \cref{ex:l07-timing} で導いたものである．}
  \label{fig:l07-timing}
\end{figure}

\begin{example}
  \label{ex:l07-timing}
  \cref{tab:l07-timing} はこれらの規則を実行例に適用したものであり，
  \cref{fig:l07-timing} はその結果を図示する．I3，I4，I6 にすべての規則
  が現れている．I3 は発行の 1 サイクル後のサイクル~4 にディスパッチされ，
  サイクル~5 に実行を終える．結果書き込みはサイクル~6 になるはずだが，
  そのときバスは I1 の結果を運ぶので，$W_3 = 7$ となる．したがって，
  サイクル~7 に加減算器が I2 のために空いているかどうかは，後の命令 I3 の
  行に依存する．I4 はサイクル~5 から準備完了だが，加減算器は，結果がバス
  を待っている I3 にその結果書き込みのサイクル~7 まで占有され，続いて
  より古い I2 にサイクル~9 まで占有されるので，$D_4 = 9$，$W_4 = 11$ と
  なる．I6 は，$W_1 = 6$ で ML1 が解放されるまで乗除算用のステーションが
  2つとも使用中なので，$I_6 = 7$ となる．そのオペランドは I2 と I4 から
  来るので，\cref{eq:l07-dispatch} より
  $D_6 \ge 1 + \max(7, 9, 11) = 12$ であり，これは I5 が乗除算器を解放
  するサイクルでもあるので，$D_6 = 12$，$W_6 = 16$ となる．この表は
  \cref{tab:l07-trace} の完全なトレースと一致する．\tip{バスが 1本なら，
  2つの命令が同じサイクルにブロードキャストすることはない．$W_n$ の列に同じ
  値が 2つあれば，バスの競合を見落としている．}
\end{example}

\begin{table}[htbp]
  \centering\small
  \caption{実行例の発行，ディスパッチ，結果書き込みのサイクル．}
  \label{tab:l07-timing}
  \begin{tabular}{@{}llccc>{\raggedright\arraybackslash}p{0.44\linewidth}@{}}
    \toprule
    & 命令 & $I_n$ & $D_n$ & $W_n$ & タイミングを決める要因 \\
    \midrule
    I1 & \texttt{MUL R1, R2, R3} & 1 & 2  & 6  & 発行時に準備完了 \\
    I2 & \texttt{ADD R4, R1, R5} & 2 & 7  & 9  & $W_1 = 6$；$W_3 = 7$ まで加減算器は I3 が使用 \\
    I3 & \texttt{SUB R6, R2, R5} & 3 & 4  & 7  & 発行時に準備完了；サイクル 6 はバスが使用中 \\
    I4 & \texttt{ADD R1, R5, R3} & 4 & 9  & 11 & 加減算器を I3，次いでより古い I2 が使用 \\
    I5 & \texttt{DIV R2, R6, R3} & 5 & 8  & 12 & $W_3 = 7$ \\
    I6 & \texttt{MUL R3, R4, R1} & 7 & 12 & 16 & $W_1 = 6$ までステーションなし；
                                               $W_4 = 11$；$W_5 = 12$ まで乗除算器が使用中 \\
    I7 & \texttt{SUB R5, R3, R2} & 8 & 17 & 19 & $W_6 = 16$ \\
    \bottomrule
  \end{tabular}
\end{table}

\begin{keypoint}[章のまとめ]
  Tomasulo のアルゴリズムは，レジスタリネーミングと，各命令をオペランドが
  揃い次第実行するアウトオブオーダ実行とを組み合わせる．発行は，命令を
  プログラム順に命令キューからリザベーションステーションへ移し，各
  オペランドをレジスタファイルの値または RAT のタグとして取得し，宛先の
  RAT エントリがそのステーションを指すようにする．ステーションは命令の
  準備ができるまで共通データバスから値を取り込み，その後ディスパッチが
  命令を空いている実行ユニットに送る．複数の命令が準備できていれば最も
  古いものを選ぶ．結果書き込みは値とタグをブロードキャストし，ステーション
  を解放し，結果が古くなっていない場合に限りレジスタファイルと RAT を更新
  する．1本のバスは 1 サイクルに 1つの結果を運び，遅いユニットが優先
  される．ロードとストアはプログラム順に実行される．1 サイクルに 1 ステップ
  という前提のもとでは，各命令の発行，ディスパッチ，結果書き込みの
  サイクルは，そのオペランド，利用可能なステーションとユニット，そして
  バスから決まる．Tomasulo のアルゴリズムは例外をサポートしない．次章の
  リオーダバッファがこれを加える．\margin{発行・ディスパッチ・結果書き込みは
  ここで述べたままである．第8章は，レジスタファイルをプログラム順に書き換える
  第 4 のステップ，コミットを加える．}
\end{keypoint}

\end{document}
